\documentclass[11pt, a4paper]{article}

% --- Geometry & Page Setup ---
\usepackage[a4paper, margin=2.5cm]{geometry}

% --- Font & Language Configuration (Compatibile con pdfLaTeX) ---
\usepackage[T1]{fontenc}
\usepackage[utf8]{inputenc}
\usepackage[english]{babel}
\usepackage{helvet}
\renewcommand{\familydefault}{\sfdefault}

% --- Essential Packages ---
\usepackage{amsmath, amssymb}
\usepackage{booktabs}
\usepackage{xcolor}
\usepackage{listings}
\usepackage{enumitem}
\usepackage{hyperref}

% --- Custom Color Definitions ---
\definecolor{brandblue}{RGB}{30, 58, 138}
\definecolor{brandaccent}{RGB}{13, 148, 136}
\definecolor{codebg}{RGB}{248, 250, 252}
\definecolor{codeframe}{RGB}{226, 232, 240}
\definecolor{commentgray}{RGB}{100, 116, 139}

% --- Setup Hyperlinks ---
\hypersetup{
    colorlinks=true,
    linkcolor=brandblue,
    citecolor=brandaccent,
    urlcolor=brandblue,
    pdftitle={Assignment Report},
    pdfauthor={Alex Morgan}
}

% --- Setup Listings (Code Environment) ---
\lstset{
    backgroundcolor=\color{codebg},
    basicstyle=\ttfamily\small,
    breaklines=true,
    captionpos=b,
    commentstyle=\color{commentgray},
    frame=single,
    rulecolor=\color{codeframe},
    keywordstyle=\color{brandblue}\bfseries,
    stringstyle=\color{brandaccent},
    showstringspaces=false,
    tabsize=2
}

% --- Custom Section Title Formatting ---
\makeatletter
\renewcommand{\section}{\@startsection{section}{1}{\z@}%
    {-3.5ex \@plus -1ex \@minus -.2ex}%
    {2.3ex \@plus .2ex}%
    {\normalfont\Large\bfseries\color{brandblue}}}
\renewcommand{\subsection}{\@startsection{subsection}{2}{\z@}%
    {-3.25ex \@plus -1ex \@minus -.2ex}%
    {1.5ex \@plus .2ex}%
    {\normalfont\large\bfseries\color{brandaccent}}}
\makeatother

% --- Document Metadata ---
\title{
    \vspace{-1cm}
    \color{brandblue}\textbf{Advanced Control Systems}\\[0.3em]
    \Large\textbf{Assignment 1: DH-Table, Direct Kinematics, Inverse Kinematics, Geometric and Analytical Jacobian}
}
\author{\textbf{Student Name:} Filippo Segalla \quad|\quad \textbf{Student ID:} VR538212\\[0.2em]
\small  Department of Engineering for Innovative Medicine (DIMI) }
\date{\today}

\begin{document}

\maketitle
\thispagestyle{empty}

\begin{abstract}
\noindent 
\end{abstract}

\hrule
\vspace{1.5em}

\section{DH-Table}
The Denavit-Hartenberg (D-H) parameter table is a standardized mathematical convention used in robotics to model the forward kinematics of serial-link manipulators. It systematically describes the relative spatial transformation between consecutive coordinate frames attached to each link along a kinematic chain using four scalar parameters: link length ($a_i$), link twist ($\alpha_i$), link offset ($d_i$), and joint angle ($\theta_i$). 

The D-H convention is widely adopted because it significantly simplifies kinematic modeling and computation. By enforcing specific geometric rules for frame placement---such as aligning the $x_i$-axis with the common normal between adjacent joint axes---the general six degrees of freedom required for 3D spatial transformations are reduced to just four parameters per joint. This structured approach allows for the systematic construction of individual homogeneous transformation matrices $\mathbf{A}_i$:

\begin{equation}
    \mathbf{A}_i = \text{Rot}_{z, \theta_i} \cdot \text{Trans}_{z, d_i} \cdot \text{Trans}_{x, a_i} \cdot \text{Rot}_{x, \alpha_i}
\end{equation}

% Passaggi di moltiplicazione (allineati sul simbolo =)
\begin{align*}
\mathbf{A}_i &=\underbrace{
\begin{bmatrix} 
\cos\theta_i & -\sin\theta_i & 0 & 0 \\ 
\sin\theta_i & \cos\theta_i & 0 & 0 \\ 
0 & 0 & 1 & 0 \\ 
0 & 0 & 0 & 1 
\end{bmatrix}}_{\text{Rot}_{z, \theta_i}}
\underbrace{
\begin{bmatrix} 
1 & 0 & 0 & 0 \\ 
0 & 1 & 0 & 0 \\ 
0 & 0 & 1 & d_i \\ 
0 & 0 & 0 & 1 
\end{bmatrix}}_{\text{Trans}_{z, d_i}}
\underbrace{
\begin{bmatrix} 
1 & 0 & 0 & a_i \\ 
0 & 1 & 0 & 0 \\ 
0 & 0 & 1 & 0 \\ 
0 & 0 & 0 & 1 
\end{bmatrix}}_{\text{Trans}_{x, a_i}}
\underbrace{
\begin{bmatrix} 
1 & 0 & 0 & 0 \\ 
0 & \cos\alpha_i & -\sin\alpha_i & 0 \\ 
0 & \sin\alpha_i & \cos\alpha_i & 0 \\ 
0 & 0 & 0 & 1 
\end{bmatrix}}_{\text{Rot}_{x, \alpha_i}}
\end{align*}

At the end of this multiplications the final DH-Matrix will exactly like this:
\begin{equation}
    \mathbf{A}_i = \begin{bmatrix} \cos\theta_i & -\sin\theta_i \cos\alpha_i & \sin\theta_i \sin\alpha_i & a_i \cos\theta_i \\ \sin\theta_i & \cos\theta_i \cos\alpha_i & -\cos\theta_i \sin\alpha_i & a_i \sin\theta_i \\ 0 & \sin\alpha_i & \cos\alpha_i & d_i \\ 0 & 0 & 0 & 1 \end{bmatrix}
\end{equation}

\noindent Concatenating these matrices yields the overall transformation matrix $\mathbf{T}_0^n$, defining the exact pose of the end-effector relative to the base coordinate system:

\begin{equation}
    \mathbf{T}_n^0 = \mathbf{A}_1(q_1) \mathbf{A}_2(q_2) \cdots \mathbf{A}_n(q_n) = \prod_{i=1}^{n} \mathbf{A}_i(q_i)
\end{equation}

\subsection{Extended Kinematic Chain: Base and Tool Transformations}

Standard Denavit-Hartenberg parameters only compute the transformation from Frame $0$ ($F_0$) to Frame $n$ ($F_n$). In real applications, we need to relate the actual tool tip to the global workspace:
\begin{itemize}[leftmargin=1.5em]
    \item \textbf{Base Frame ($F_B$):} The global reference system (e.g., workbench or room origin).
    \item \textbf{End-Effector Frame ($F_E$):} The frame attached to the Tool Center Point (TCP).
\end{itemize}

To compute the complete position of the TCP relative to the world, we multiply the joint transformation matrix $\mathbf{T}_n^0(q)$ by two constant matrices:

\begin{equation}
    \mathbf{T}_E^B(q) = \mathbf{T}_0^B \cdot \mathbf{T}_n^0(q) \cdot \mathbf{T}_E^n = \mathbf{T}_0^B \cdot \left( \prod_{i=1}^{n} \mathbf{A}_i(q_i) \right) \cdot \mathbf{T}_E^n
\end{equation}

% --- Base Frame Transformation ---
\subsubsection{Base Frame Transformation \texorpdfstring{($\mathbf{T}_0^B$)}{(T0B)}}
Frame $0$ is fixed on Joint 1 and does not always sit at the world origin $F_B$. The matrix $\mathbf{T}_0^B$ offsets the robot base by position $(x_b, y_b, h_b)$ and rotation $\gamma_b$:

\begin{equation}
\mathbf{T}_0^B = 
\begin{bmatrix} 
\cos\gamma_b & -\sin\gamma_b & 0 & x_b \\ 
\sin\gamma_b &  \cos\gamma_b & 0 & y_b \\ 
0            &  0           & 1 & h_b \\ 
0            &  0           & 0 & 1 
\end{bmatrix}
\end{equation}

\noindent \textbf{Why it is used:}
\begin{itemize}[leftmargin=1.5em]
    \item $(x_b, y_b, h_b)$ is the physical 3D distance from the world origin to the robot base.
    \item $\gamma_b$ is the rotation angle of the robot mounting base.
    \item If the robot is moved to a new table, only $\mathbf{T}_0^B$ changes; internal robot parameters stay the same.
\end{itemize}

% --- End-Effector Transformation ---
\subsubsection{Tool Transformation \texorpdfstring{($\mathbf{T}_E^n$)}{(TEn)}}
Frame $n$ lies at the wrist flange, but tasks happen at the tool tip (TCP). The matrix $\mathbf{T}_E^n$ describes an extension $L_t$, a lateral shift $d_x$, and a tool rotation of $-90^\circ$:

\begin{equation}
\mathbf{T}_E^n = 
\begin{bmatrix} 
0 & 0 & -1 & d_x \\ 
0 & 1 &  0 & 0   \\ 
1 & 0 &  0 & L_t \\ 
0 & 0 &  0 & 1 
\end{bmatrix}
\end{equation}

\noindent \textbf{Why it is used:}
\begin{itemize}[leftmargin=1.5em]
    \item $(d_x, 0, L_t)$ is the physical distance from the robot wrist to the end-effector tip.
    \item The rotation aligns the tool orientation (e.g., pointing the fingers forward).
    \item When replacing the tool (e.g., gripper to torch), only $\mathbf{T}_E^n$ needs to be updated.
\end{itemize}

\end{document}